%%
%% %%%

%%%   Самарский государственный технический университет

%%%   Кафедра прикладной математики и информатики

%%%

%%%   Преамбула для статьи в журнал "Вестник Самарского государственного

%%%   технического университета. Серия: «Физико-математические науки»"

%%%   Ориентирована под MikTEx 2.4 for Win32

%%%

%%%   Хотя сам журнал собирается под GNU/Linux на дистрибутиве TeXLive



\documentclass[11pt,a4paper,twoside]{article}



\usepackage[cp1251]{inputenc}

\usepackage[T2A]{fontenc}

\usepackage[english,russian]{babel}

\usepackage{samgtubib}

\usepackage[dvips]{graphicx}

\usepackage[multiple]{footmisc}



\usepackage{samgtu}

\renewcommand{\VestnikYear}{2009} % Год издания Вестника

\renewcommand{\VestnikNumber}{2\,(19)} % Номер Вестника

\renewcommand{\VestnikMonth}{Ноябрь} % Месяц выхода Вестника

\renewcommand{\VestnikData}{01.11.09} % Дата подписания Вестника в печать

\renewcommand{\VestnikUPL}{26,64} % Усл. печ. л.

\renewcommand{\VestnikUIL}{25,73} % Уч.-изд. л.

\renewcommand{\VestnikRegNumber}{479/09} % Регистрационный номер

\renewcommand{\VestnikOrder}{994} % Номер заказа






%% 
%%% \usepackage{itemize}
%%%
% \begin{document}


\renewcommand{\firstpage}{26}
\renewcommand{\lastpage}{33}
% \newcommand{\bigcom}{\text{\usefont{T2A}{cmfib}{m}{n}\large{\raisebox{1pt}{,}}}} % Cимвол производной в ТФКП (ковариантная)
% \Partition{Дифференциальные уравнения}{Applied mechanics} % Раздел Вестника, в который подаётся статья
% Названия разделов см. на сайте при подаче заявки


%%%%%%%%%%%%%%%%%%%%%%%%%%%%%%%%%%%%%%%%%%%%%%%%%%%%%%%%%%%%%%%%%%%%%%%%%%%%%%%%%%%%
%Информация о статье и об авторах на русском и английском языках

% Номер УДК
\udk{517.91/.93}
\msc{34G10; 47N20}% Коды MSC см. http://www.ams.org/msc/
\title{Уравнения соболевского типа второго порядка с~относительно диссипативным пучком операторов}% Полный заголовок
{Уравнения соболевского типа второго порядка \dots}% Заголовок, который идёт в верхний колонтитул

\author{А.\,А.~Замышляева, О.\,Н.~Цыпленкова}% Автор(ы) для заголовка, если авторов несколько укать \inst
{З\,а\,м\,ы\,ш\,л\,я\,е\,в\,а~А.\,А.,\ Ц\,ы\,п\,л\,е\,н\,к\,о\,в\,а~О.\,Н.}% Автор(ы) для оглавления и колонтитула через \,

\institute{\yurgu.}

\email{E-mails}{alzama@mail.ru, Tsyplenkova\_Olga@mail.ru} %E-mail's авторов

\otherinf{\fullauthor{Алёна Александровна Замышляева}\regalia{к.ф.-м.н., доц.}\position{докторант}
\dept{каф. уравнений математической физики} 
\fullauthor{Ольга Николаева Цыпленкова} \position{аспирант} \dept{каф. уравнений математической физики} }

\keywords{уравнение соболевского типа, относительная диссипативность пучка операторов, аккретивные операторы, фазовое пространство}

\workabstract{Рассмотрена задача Коши для уравнения соболевского типа второго порядка. 
Введено определение относительной диссипативности пучка операторов, обобщено понятие диссипативности и~относительной диссипативности оператора.
Установлена связь с~результатами теории аккретивных операторов.
Согласно идеологии Келдыша, исходная задача редуцируется к~задаче Коши для уравнения соболевского типа первого порядка. Приводятся результаты для исследуемой задачи.}

\engtitle{The Sobolev-type equations of the second order with~the~relatively dissipative operator pencils}

\engauthor{A.\,A.~Zamyshlyaeva, O.\,N.~Tsyplenkova}{Z\,a\,m\,y\,s\,h\,l\,y\,a\,e\,v\,a~A.\,A.,
T\,s\,y\,p\,l\,e\,n\,k\,o\,v\,a~O.\,N.}


\enginstitute{\engyurgu.}

\otherenginfm{\fullauthor{Alyona A. Zamyshlyaeva} \regalia{Ph.\,D. (Phys. \& Math.)} \position{Doctoral Candidate} \dept{Dept. of Mathematical Physics Equations} \fullauthor{Ol'ga N. Tsyplenkova} \position{Postgraduate Student} \dept{Dept. of Mathematical Physics Equations}}

\engabstract{Of concern is the Cauchy problem for the Sobolev-type equation of the second order. We introduce the definition of relatively dissipative operator pencils, generalize the notion of dissipativity and relative dissipativity of operators. The connection with the theory of accretive operators is established. According to the Keldysh ideology, the original problem is reduced to the Cauchy problem for the Sobolev-type equation of the first order and the results for the investigated problem are obtained. }

\engkeywords{Sobolev-type equation, relative dissipativity of the operator pencil, accretive operators, phase space}

\date{19/X/2011} % Дата поступления статьи в редакцию.
\revisiondate{12/II/2012} % В окончательном виде

%%%%%%%%%%%%%%%%%%%%%%%%%%%%%%%%%%%%%%%%%%%%%%%%%%%%%%%%%%%%%%%%%%%%%%%%%%%%%%%%%%%%


\maketitle


\Section[n]{Введение} 
Пусть $\mathcal V$ и $\mathcal G$ "--- гильбертовы пространства со скалярными произведениями \mbox{$\langle{}\cdot{}, {}\cdot{}\rangle$} и
$[{}\cdot{}, {}\cdot{}]$ соответственно, оператор $A \in \mathcal L
(\mathcal V;\mathcal G)$, а операторы $B_{1},$ $B_{0} \in
\mathcal Cl (\mathcal V;\mathcal G)$ (линейны, замкнуты, плотно
определены в $\mathcal V$).

Рассмотрим задачу Коши
\begin{equation}
v(0)=v_{0}, \quad \dot{v}(0)=v_{1} 
\label{tsyp:1}
\end{equation}
для операторно"=дифференциального уравнения 
\begin{equation}
A\ddot{v}=B_{1}\dot{v}+B_{0}v.
\label{tsyp:2}
\end{equation} 

Данная задача рассматривалась многими авторами, в частности в [\citen{zamtsyp:bib:Zamyshlyaeva}--\citen{zamtsyp:bib:Showalter}]. R.~Showalter исследовал задачу \eqref{tsyp:1}, \eqref{tsyp:2} в случае непрерывной обратимости оператора $A$, когда оператор
($-A^{-1}B_{1}$) непрерывен и аккретивен, а оператор ($-A^{-1}B_{0}$) регулярно {$m$"=аккретивен.} 
В~работе A.~Favini и A.~Yagi накладываются условия самосопряжённости и неотрицательности на оператор~$A$, самосопряжённости и положительности на оператор $B_{0}$ и диссипативности на оператор $B_{1}$. 
Нас интересует разрешимость задачи \eqref{tsyp:1}, \eqref{tsyp:2} в случае необратимости оператора~$A$.

Неполное уравнение соболевского типа
\begin{equation}
A\ddot{v}=Bv, 
\label{tsyp:3}
\end{equation}
где операторы $A,$ $ B\in {\mathcal L} (\mathcal V,\mathcal G)$,
причём оператор $B$ $(A,p)$-ограничен, было впервые изучено в~\cite{zamtsyp:bib:Sviridyuk}.
Единственное решение $v\in C^\infty(\Bbb R;\mathcal V)$ задачи Коши \eqref{tsyp:1} для уравнения \eqref{tsyp:3} представимо в виде $$v(t)=V^{t}_{1}v_1+V^{t}_{0}v_0, $$
где \textit{пропагаторы} $V^{t}_{k}$ имеют следующий вид:
 $$V^{t}_{k}=\frac{1}{2 \pi i}\int_\gamma\mu^{1-k}\left(\mu^2A-B\right)^{-1}Ae^{\mu t}d\mu,\quad k=0, 1. $$ 
Здесь контур $\gamma\subset \Bbb C$ ограничивает область, содержащую $A$-спектр оператора~$B$, а начальные значения
$v_k \in \mathop{\rm im} V^{0}_{1}=\mathop{\rm im} V^{0}_{0}$, $k=0, 1$,
где $\mathop{\rm im} \dot{V}^{0}_{1}=\mathop{\rm im} V^{0}_{0}$ "---
подпространство в $\mathcal{V}$.

Полное уравнение соболевского типа второго порядка \eqref{tsyp:2} было
изучено ранее \cite{zamtsyp:bib:Zamyshlyaeva} в~случае относительно полиномиальной
ограниченности пучка операторов $B_{0}$ и $B_{1}$ (в~дальнейшем пучок будем обозначать через $\vec{B}$). 
Здесь построено аналитическое семейство $M$, $N$-функций и фазовое пространство данного уравнения.

Согласно идеологии М.\,В.~Келдыша, в данной работе уравнение \eqref{tsyp:2}
редуцируется к~эквивалентному ему уравнению соболевского типа первого порядка. 
Наш подход заключается в~построении фазового пространства такого уравнения.

Работа содержит 4 пункта. 
В~первом даны основные определения и~результаты теории уравнений соболевского типа, необходимые в~дальнейших исследованиях. 
Во втором пункте приводятся результаты, касающиеся $L$"=диссипативных операторов из \cite{zamtsyp:bib:Fedorov}, и вводится понятие $A$"=диссипативности пучка операторов $\vec{B}$. 
Третий пункт содержит некоторые результаты теории аккретивных операторов \cite{zamtsyp:bib:Showalter}. 
Здесь же доказывается, что результаты данной работы являются более общими по сравнению с~\cite{zamtsyp:bib:Showalter}. 
Основной результат "--- теорема о~существовании и~единственности решения задачи \eqref{tsyp:1}, \eqref{tsyp:2} "--- приведён в~четвёртом пункте работы.

\smallskip

\Section{Уравнения соболевского типа с относительно радиальными операторами}
Пусть $\frak U$ и $\frak F$ "--- гильбертовы пространства, оператор $L\in \mathcal L (\frak U;\frak F)$, а оператор  $M: \mathop{\rm dom} M\subset\frak U\rightarrow\frak F$ линеен и замкнут, плотно определён.

\smallskip

\begin{definition}[1.1]Множество $$\rho^{L}(M)=\{\mu\in\mathbb{C}:(\mu L-M)^{-1}\in\mathcal L
(\frak F;\frak U)\}$$ называется \emph{резольвентным множеством} оператора $M$ \emph{относительно} оператора $L$ (короче, \emph{$L$-резольвентным множеством} оператора $M$). 
Множество $\mathbb{C}\backslash \rho^{L}(M)=\sigma^{L}(M)$ называется \emph{спектром} оператора $M$ \emph{относительно} оператора $L$ (короче, \emph{$L$-спектром} оператора $M$).
\end{definition}

\smallskip

\begin{definition}[1.2]Оператор-функции $(\mu L-M)^{-1}$,
$R^{L}_{\mu}(M)=(\mu L-M)^{-1}L$, $L^{L}_{\mu}(M)=L(\mu L-M)^{-1}$  с областью определения $\rho^{L}(M)$ называются
соответственно \emph{резольвентой, правой резольвентой, левой
резольвентой} оператора \emph{$M$ относительно} оператора $L$
(короче, \emph{$L$"=резольвентой, правой  $L$"=резольвентой, левой
$L$"=резольвентой оператора $M$}).
\end{definition}

\smallskip

\begin{definition}[1.3]Оператор $M$ называется \emph{$(L, 0)$-радиальным относительно} оператора $L$, если
\begin{itemize}
\item[(i)] $\exists a\in\mathbb{R}$ $\forall\mu>a$ $\mu\in\rho^{L}(M)$;
\item[(ii)] $\exists K\in \mathbb{R}_{+}$ $\forall\mu>a$ $\forall n\in\mathbb{N}$
$$\max\bigl\{ \| (R^{L}(M))^{n} \|  _{\mathcal{L}(\mathcal{U})},
 \| (L^{L}(M))^{n} \|  _{\mathcal{L}(\mathcal{F})}\bigr\}\leq
\frac{K}{ (\mu -a)^{n}}.$$
\end{itemize}
\end{definition}

\smallskip

\begin{remark}[1.1] Без потери общности в определении 1.3 можно
положить \mbox{$a=0$.}  И в дальнейшем будем считать, что $a=0$.
 \end{remark}

\smallskip

\begin{newthm}{Утверждение 1.1} Пусть оператор $L$ непрерывно обратим{\rm .} 
Тогда оператор $L^{-1}M\in\mathcal Cl(\frak{U})$ $($или $ML^{-1})$ радиален точно тогда{\rm ,} когда оператор $M$
$(L, 0)$"=радиален{\rm .}
\end{newthm}

\smallskip

\begin{definition}[1.4]Множество $\mathcal{P}\subset\mathcal{D}$
называется \emph{фазовым пространством}  уравнения 
\begin{equation}
L\dot{u}=Mu, 
\label{tsyp:4}
\end{equation} 
если
\begin{itemize}
\item[(i)] любое решение $u=u(t)$ уравнения \eqref{tsyp:4} лежит в $\mathcal{P}$,
т.~е. $u(t)\in\mathcal{P}$ $\forall t\in [0,T]$, $T>0$;
\item[(ii)] для любого $u_{0}$ из некоторого плотного в $\mathcal{P}$
множества существует единственное решение задачи
\begin{equation}
u(0)=u_0 
\label{tsyp:5}
\end{equation}
для уравнения \eqref{tsyp:4}.
\end{itemize}
\end{definition}

\smallskip


\begin{theorem}[1.1 \cite{zamtsyp:bib:Fedorov1}]Пусть оператор $M$ является $(L, 0)$-радиальным{\rm .}
Тогда \mbox{$\frak U^{1}= \overline{im(\mu L-M)^{-1}L}$} есть фазовое пространство уравнения \eqref{tsyp:4}{\rm .} 
\end{theorem}

\smallskip

\Section{$\boldsymbol{L}$-диссипативные операторы}
 Пусть $\frak U$ и $\frak F$ "--- гильбертовы  пространства, оператор
$L\in \mathcal L (\frak U;\frak F)$, оператор $M\in \mathcal
Cl (\frak U;\frak F)$. Рассмотрим задачу Коши \eqref{tsyp:5} для
операторно"=дифференциального уравнения \eqref{tsyp:4}.

\smallskip

\begin{definition}[2.1]Оператор $M$ будем называть \emph{диссипативным относительно} оператора $L$ (короче, \emph{$L$"=диссипативным}) по отношению к скалярным произведениям $\langle{}\cdot{}, {}\cdot{}\rangle$ в $\frak U$ и~$[{}\cdot{}, {}\cdot{}]$ в $\frak F$, если
\begin{itemize}
\item[(i)] $\forall u\in \mathop{\rm \mathop{\rm dom}}  M $ $\mathop{\rm Re}[ L u, M u]\leq 0$; 
\item[(ii)] существует положительное число $\alpha \in \rho^{L}(M) $;
\item[(iii)] $\forall u\in \mathop{\rm \mathop{\rm dom}} M $ $\mathop{\rm Re}\langle(\alpha L - M )^{-1} L u,(\alpha L - M )^{-1} M u\rangle\leq 0$.
\end{itemize}
\end{definition}

\smallskip

\begin{remark}[2.1]  Если условие (i) определения 2.1 записать в~эквивалентном виде
$$\forall f\in \frak F \quad \mathop{\rm Re}[L(\alpha L-M)^{-1}f,M(\alpha
L-M)^{-1}f]\leq 0,$$ то условия (i) и (iii) примут симметричный
вид.
\end{remark}

\smallskip

\begin{lemma}[2.1]Пусть выполняется условие {\rm(i)} определения
{\rm2.1,} тогда условие {\rm (ii)} эквивалентно следующему\/{\rm :} $$(\mathop{\rm ker} L\cap \mathop{\rm ker}
M=\{0\})\wedge (im(\alpha L-M)= \frak F).$$
\end{lemma}

\smallskip

\begin{theorem}[2.1]Если оператор $M$ $L$-диссипативен{\rm ,} тогда он {$(L, 0)$"=радиален{\rm .}}
\end{theorem}

\smallskip

Пусть $\mathcal V$ и $\mathcal G$ "--- гильбертовы пространства со скалярными произведениями $\langle{}\cdot{}, {}\cdot{}\rangle$ и $[{}\cdot{}, {}\cdot{}]$ соответственно, 
операторы $A \in \mathcal L (\mathcal V;\mathcal G)$, а операторы \mbox{$B_{1},B_{0} \in \mathcal Cl (\mathcal V;\mathcal G)$.}

Рассмотрим задачу Коши \eqref{tsyp:1} для операторно"=дифференциального уравнения \eqref{tsyp:2}.

\smallskip

\begin{definition}[2.2]Операторный пучок $\vec{B}$ будем называть \emph{диссипативным относительно} оператора $A$ (короче,
\emph{$A$"=диссипативным}\/)  по отношению к скалярным произведениям $\langle{}\cdot{}, {}\cdot{}\rangle$ в $\cal V$ и $[{}\cdot{}, {}\cdot{}]$ в $\cal G$, если
\begin{itemize}
\item[(i)] $\forall v_{1},v_{2}\in \mathop{\rm \mathop{\rm dom}}\vec{B} (\mathop{\rm \mathop{\rm dom}}\vec{B}=\mathop{\rm \mathop{\rm dom}}B_{1}\bigcap
\mathop{\rm \mathop{\rm dom}}B_{0})$ $\mathop{\rm Re}(\langle
v_{1},v_{2}\rangle)+[Av_{2},B_{0}v_{1}+B_{1}v_{2}]\leq0$;
\item[(ii)] существует положительное число $\alpha \in
\rho^{A}(\bar{B})$;
\item[(iii)] $\forall v_{1},v_{2} \in  \mathop{\rm dom}\vec{B}$ $\mathop{\rm Re} \langle
R_{\alpha} B_{0} v_{1}+\alpha R_{\alpha} A v_{2},R_{\alpha}
(\alpha A-B_{1}+\alpha B_{0}) v_{1}+R_{\alpha} ( A+\alpha B_{1}+
B_{0})v_{2} \rangle \leq 0$, где  $ R_{\alpha}=(\alpha^{2}
A-\alpha B_{1}-B_{0})^{-1}$ "--- относительная
$A$"=резольвента \emph{пучка $\vec{B}$}.
\end{itemize}
\end{definition}

\smallskip

\begin{lemma}[2.2]Пусть выполнено условие {\rm (i)} определения {\rm 2.2,}
тогда условие {\rm (ii)} эквивалентно следующему\/{\rm:} $$(\mathop{\rm ker} A\cap \mathop{\rm ker}
B_{1}\cap \mathop{\rm ker} B_{0}=\{0\})\wedge (im(\alpha^{2}A-\alpha
B_{1}-B_{0})=\mathcal G).$$ 
\end{lemma}

\smallskip 

\begin{proof}
Покажем, что из (i) и тривиальности пересечения ядер $\mathop{\rm ker} A$, $\mathop{\rm ker}
B_{1}$ и $\mathop{\rm ker} B_{0}$ следует инъективность оператора $\alpha^{2} A-\alpha
 B_{1}-B_{0}$, $\alpha >0.$ 
 Пусть $\alpha^{2} A=\alpha  B_{1}+B_{0}$ и возьмём $v_{2}=\alpha u$, $v_{1}=u$. Тогда
\begin{multline*}
\mathop{\rm Re}(\langle u,\alpha u\rangle+[A\alpha
u,B_{0}u+B_{1}\alpha u])=\mathop{\rm Re}(\alpha\langle u,
u\rangle+\alpha[A u,\alpha^{2} A u])=
\\
=\mathop{\rm Re}(\alpha\langle u, u\rangle+\alpha^{3}[A u, A u])=
\mathop{\rm Re}(\alpha\|u\|^{2}+\alpha^{3}\|A u\|^{2})\leq 0.
\end{multline*}
Отсюда  следует, что $u\in (\mathop{\rm ker} A\cap \mathop{\rm ker} B_{1}\cap \mathop{\rm ker} B_{0})$,
то есть $u=0.$ 
\end{proof}

\smallskip


Сведём задачу \eqref{tsyp:1}, \eqref{tsyp:2} к задаче \eqref{tsyp:4}, \eqref{tsyp:5}, где
$u(t)=\begin{pmatrix} v(t)\\v'(t)\end{pmatrix},$ операторы
\mbox{$L=\begin{pmatrix} I&0\\0&A\end{pmatrix}\in\mathcal L (\frak
U;\frak F),$} $M=\begin{pmatrix}
0&I\\B_{0}&B_{1}\end{pmatrix}\in\mathcal Cl (\frak U;\frak F),$
пространства $\frak U=\mathcal V\times \mathcal V,$ $\frak
F=\mathcal V\times \mathcal G.$

\smallskip

\begin{lemma}[2.3]Пучок операторов $\vec{B}$ является $A$"=диссипативным точно тогда{\rm ,} когда оператор $M$ является $L$"=диссипативным{\rm .}
\end{lemma}

\smallskip

\begin{proof} Докажем \emph{необходимость}.

\begin{itemize}

\item[(i)] Докажем выполнение условия (i) в определении 2.2:
\begin{multline*}
\mathop{\rm Re}(\langle v_{1},v_{2}\rangle_{\mathcal
V}+[Av_{2},B_{0}v_{1}+B_{1}v_{2}]_{\mathcal G})=
\mathop{\rm Re}\left(\begin{pmatrix}
v_{1}\\Av_{2}\end{pmatrix},\begin{pmatrix}
v_{2}\\B_{0}v_{1}+B_{1}v_{2}\end{pmatrix}\right)=
\\
=\mathop{\rm Re}\left(\begin{pmatrix}
I&0\\0&A\end{pmatrix}\begin{pmatrix}v_{1}\\v_{2}\end{pmatrix}
,\begin{pmatrix}
0&I\\B_{0}&B_{1}\end{pmatrix}\begin{pmatrix}v_{1}\\v_{2}\end{pmatrix}\right)=\mathop{\rm Re}[Lu,Mu]\leq 0.
\end{multline*}

\item[(ii)] Пусть существует положительное число $\mu\in\rho^{L}(M)$,
тогда существует $(\mu
L-M)^{-1}\in\mathcal{L}(\frak{F},\frak{U})$, при этом
\begin{multline*}
(\mu L-M)^{-1}=\left(\mu \begin{pmatrix}
I&0\\0&A\end{pmatrix}-\begin{pmatrix}
0&I\\B_{0}&B_{1}\end{pmatrix}\right)^{-1}= \\ = \left(
\begin{pmatrix} \mu&0\\0&\mu A\end{pmatrix}-\begin{pmatrix}
0&I\\B_{0}&B_{1}\end{pmatrix}\right)^{-1}=
\\
=\begin{pmatrix} \mu&-I\\-B_{0}&\mu
A-B_{1}\end{pmatrix}^{-1}=(\mu^{2}A-\mu
B_{1}-B_{0})^{-1}\begin{pmatrix} \mu
A-B_{1}&I\\B_{0}&\mu\end{pmatrix}=
\\
=\begin{pmatrix} (\mu A-B_{1})(\mu^{2}A-\mu
B_{1}-B_{0})^{-1}&(\mu^{2}A-\mu
B_{1}-B_{0})^{-1}\\B_{0}(\mu^{2}A-\mu B_{1}-B_{0})^{-1}&\mu
(\mu^{2}A-\mu B_{1}-B_{0})^{-1}\end{pmatrix}.
\end{multline*}

Отсюда следует существование $(\mu^{2}A-\mu
B_{1}-B_{0})^{-1}\in\mathcal L (\mathcal G;\mathcal
 V)$, то есть $\mu\in\rho^{A}(\vec{B})$, где $\rho^{A}(\vec{B})=\{\mu\in\mathbb{C}:(\mu^{2}A-\mu B_{1}-B_{0})^{-1}\in\mathcal{L}(\mathcal{G},\mathcal{V})\}$.

\item[(iii)] Докажем выполнение пункта (iii) в определении 2.2:
\begin{multline*}
\mathop{\rm Re}(\langle R_{\alpha} B_{0}v_{1}+\alpha R_{\alpha}Av_{2},R_{\alpha}(\alpha A-B_{1}+\alpha B_{0})v_{1}+R_{\alpha}(A+\alpha B_{1}+B_{0})v_{2}\rangle)=
\\
=\mathop{\rm Re}\left(R_{\alpha}\begin{pmatrix} \alpha A-B_{1}&A\\B_{0}&\alpha A\end{pmatrix}\begin{pmatrix} v_{1}\\v_{2}\end{pmatrix}, R_{\alpha}\begin{pmatrix} B_{0}&\alpha A\\\alpha B_{0}&B_{0}+\alpha B_{1}\end{pmatrix}\begin{pmatrix} v_{1}\\v_{2}\end{pmatrix}\right)=
\\
=\mathop{\rm Re}\left(R_{\alpha}\begin{pmatrix} \alpha A-B_{1}&I\\B_{0}&\alpha \end{pmatrix}L\begin{pmatrix} v_{1}\\v_{2}\end{pmatrix},R_{\alpha}\begin{pmatrix}\alpha A- B_{1}&I\\ B_{0}&\alpha\end{pmatrix}M\begin{pmatrix} v_{1}\\v_{2}\end{pmatrix}\right)=
\\
=\mathop{\rm Re}\left(\begin{pmatrix} \alpha &-I\\-B_{0}&\alpha A-B_{1} \end{pmatrix}^{-1}L\begin{pmatrix} v_{1}\\v_{2}\end{pmatrix}, \begin{pmatrix}\alpha&-I\\ -B_{0}&\alpha A-B_{1}\end{pmatrix}^{-1}M\begin{pmatrix} v_{1}\\v_{2}\end{pmatrix}\right)=
\\
=\mathop{\rm Re}\left(\left(\alpha\begin{pmatrix} I &0\\0&A \end{pmatrix}-\begin{pmatrix} 0&I\\B_{0}&B_{1}\end{pmatrix}\right)^{-1}\begin{pmatrix} I&0\\0&A\end{pmatrix}\begin{pmatrix} v_{1}\\v_{2}\end{pmatrix},(\alpha L-M)^{-1}Mu\right)=
\\
=\mathop{\rm Re}((\alpha L-M)^{-1}Lu,(\alpha L-M)^{-1}Mu)\leq 0.
\end{multline*}


\end{itemize}

Доказательство  \emph{достаточности} проводится аналогично.
\end{proof}

\smallskip

\Section{Аккретивные операторы}
 Пусть $H$ "--- гильбертово пространство, $D$ "--- подпространство в
$H$ и пусть оператор {$A:D\rightarrow H$} линеен.

\smallskip

\begin{definition}[3.1]Неограниченный оператор $A:D\rightarrow H$ называется \emph{аккретивным},  если
$$(Au,u)_{H}\geq 0, \quad  x\in D,$$ и \emph{$m$"=аккретивным}, если он аккретивен и $\mathop{\rm Im}(A+I)=H.$
\end{definition}

\smallskip

\begin{newthm}{Утверждение 3.1} Следующие условия эквивалентны\/{\rm :}
\begin{itemize}
\item[\rm (i)] оператор $A:D\rightarrow H$ аккретивен и существует
$\mu\geq 0${\rm ,} что $\mathop{\rm Im}(\mu I+A)=H$\/{\rm ;}
\item[\rm (ii)] оператор $A$ "--- $m$-аккретивен\/{\rm ;}
\item[\rm (iii)] оператор $A$ аккретивен с областью определения $D,$
плотной в $H${\rm ,} и $\mathop{\rm Im}(\lambda I+A)=H$ для всех $\lambda > 0.$
\end{itemize}
\end{newthm}

\smallskip

\begin{definition}[3.2]Неограниченный оператор $A:D\rightarrow H$ называется \emph{регулярно \mbox{$m$"=аккретивным,}} если для всех $\varepsilon>0$ $\mathop{\rm Im}(A+\varepsilon I)=H$.
\end{definition}

\smallskip

Пусть $V=G=H.$

\smallskip

\begin{remark}[3.1] Пусть существует оператор $A^{-1}\in \mathcal L (\mathcal G;\mathcal V).$ Если оператор $(-A^{-1}B_{1})$
непрерывен и аккретивен, а оператор ($-A^{-1}B_{0}$) регулярно $m$"=аккретивен, то пучок $\vec{B}$ является \mbox{$A$-диссипативным.}
\end{remark}

\textit{Д\,о\,к\,а\,з\,а\,т\,е\,л\,ь\,с\,т\,в\,о.} Без ограничения общности положим $A=I$.

\begin{itemize}
\item[(i)] $\mathop{\rm Re}(\langle v_{1},v_{2}\rangle+[v_{2},B_{0}v_{1}+B_{1}v_{2}])=
\mathop{\rm Re}((\langle v_{1},v_{2}\rangle+[v_{2},B_{0}v_{1}]+[v_{2},B_{1}v_{2}])=
\mathop{\rm Re}((\langle v_{1},v_{2}\rangle+[B_{0}v_{1},v_{2}]+[v_{2},B_{1}v_{2}])$.

Рассмотрим скалярное произведение в $\frak U=H\times H$:
$$(\bar{u},\bar{v})_{\frak U}=(-B_{0}u_{1},v_{1})+(u_{2},v_{2}).$$
Оно эквивалентно естественному скалярному произведению в
$\frak{U}$, так как оператор $(-B_{0})$ регулярно $m$-аккретивный.

Операторы $L=\begin{pmatrix} I&0\\0&A\end{pmatrix}\in\mathcal L
(\frak U;\frak F),$ $M=\begin{pmatrix}
0&I\\B_{0}&B_{1}\end{pmatrix}\in\mathcal Cl (\frak U;\frak F),$
пространства $\frak U=H\times  H,$ $\frak F= H\times  H.$ Причём
\begin{multline*}
(L\bar{u},M\bar{u})_{\frak U}=  (\bar{u},M\bar{u})_{\frak U}=
\left(\begin{pmatrix}u_{1}\\u_{2}\end{pmatrix},\begin{pmatrix}u_{2}\\B_{0}u_{1}+B_{1}u_{2}\end{pmatrix}\right)_{\frak U}=
\\
=(-B_{0}u_{1},u_{2})+(u_{2},B_{0}u_{1}+B_{1}u_{2})=(B_{1}u_{2},u_{2})\leq 0.
\end{multline*}

\item[(ii)] Рассмотрим уравнение \eqref{tsyp:2}. Положим $w(t)=e^{-\alpha t}
v(t).$ Выражая $v(t)$ и подставляя в \eqref{tsyp:2}, получаем
$$\ddot{w}(t)+(2\alpha I-B_{1})\dot{w}(t)+(\alpha^{2}I-\alpha
B_{1}-B_{0})w(t)=0.$$ Оператор $B_{1}$ является аккретивным, т.~е.
$(-B_{1}u,u)_{H}\geq 0$.

Рассмотрим скалярное произведение в $\frak U=H\times H$:
$$(\bar{u},\bar{v})_{\frak U}=(-B_{0}u_{1},v_{1})+(u_{2},v_{2}).$$
Оно эквивалентно естественному скалярному произведению в
$\frak{U}$, так как оператор $-B_{0}$ регулярно $m$-аккретивный.

Определим оператор $\mathbb{A}$ следующим образом:
$\mathbb{A}\bar{u}=\begin{pmatrix}-u_{2}\\-B_{0}u_{1}-B_{1}u_{2}\end{pmatrix}$,
$\forall u_{1},u_{2}\in \mathop{\rm \mathop{\rm dom}}\vec{B}  (\mathop{\rm \mathop{\rm dom}}\vec{B}=\mathop{\rm \mathop{\rm dom}}B_{1}\bigcap
\mathop{\rm \mathop{\rm dom}}B_{0})$. Это оператор, который появляется, когда мы запишем
уравнение \eqref{tsyp:2} в виде системы:
$$\left\{\begin{array}{l}\dot{v}(t)-u(t)=0,\\\dot{u}(t)-B_{1}u(t)-B_{0}v(t)=0;
\end{array}\right. $$
в таком случае функция
$\bar{u}(t)=\begin{pmatrix}u(t)\\v(t)\end{pmatrix}$ при $t>0$
является решением уравнения
\mbox{$\dot{\bar{u}}(t)+\mathbb{A}\bar{u}(t)=\bar{0}$.}

Пусть $\bar{f}=\begin{pmatrix}f_{1}\\f_{2}\end{pmatrix}\in \frak
F$, попытаемся найти
$\bar{u}=\begin{pmatrix}u_{1}\\u_{2}\end{pmatrix}$ такой, что
$\alpha\bar{u}+\mathbb{A}\bar{u}=\bar{f}$, то есть
$$\left\{\begin{array}{l}(\alpha^{2}I-\alpha
B_{1}-B_{0})u_{1}=(\alpha-B_{1})f_{1}+f_{2},\\u_{2}=\alpha
u_{1}-f_{1}.
\end{array}\right. $$
Билинейная форма, определённая в $H$ формулой
$$a(u,v)=\alpha^{2}(u,v)_{H}+\alpha(-B_{1}u,v)_{H}+(-B_{0}u,v)_{H},$$
непрерывна и коэрцитивна. Требуемый результат следует из теоремы
Лакса"--~Мильграма: для каждого $h\in H$ существует единственное
$u\in H$ такое, что $$a(u,v)=(h,v)_{H},$$ а это равносильно
$(\alpha^{2}I-\alpha B_{1}-B_{0})u=h$. Следовательно,
$(\alpha^{2}I-\alpha B_{1}-B_{0})=H$.


\item[(iii)] Так как $A=I$, то и $L=I$. Имеем
$\langle(\alpha L-M)^{-1}L\bar{u},(\alpha L-M)^{-1}M\bar{u}\rangle_{\frak
U}=\langle(\alpha I-M)^{-1}\bar{u},(\alpha I-M)^{-1}M\bar{u}\rangle_{\frak U};$
$(\alpha I-M)^{-1}M=-(\alpha I-M)^{-1}(\alpha I-M-\alpha
I)=-I+\alpha(\alpha I-M)^{-1}=-(\alpha I-M)(\alpha
I-M)^{-1}+\alpha(\alpha I-M)^{-1}=(-\alpha I+M+\alpha I)(\alpha
I-M)^{-1}=M(\alpha I-M)^{-1}.$

Обозначив $(\alpha I-M)^{-1}\bar{u}=\bar{v}$, получаем
$\langle(\alpha I-M)^{-1}\bar{u},(\alpha I-M)^{-1}M\bar{u}\rangle_{\frak
U}=\langle(\alpha I-M)^{-1}\bar{u},M(\alpha I-M)^{-1}\bar{u}\rangle_{\frak
U}=\langle\bar{v}, M\bar{v}\rangle_{\frak U}\leq 0$ в силу (i). $\square$


\end{itemize}


\smallskip

\Section{Задача Коши для уравнения соболевского типа второго порядка с относительно диссипативным пучком операторов}

\begin{definition}[4.1]\emph{Решением задачи} \eqref{tsyp:1}, \eqref{tsyp:2} будем называть функцию $v(t)\in {\mathcal C}^2([0,T],{\mathcal V})\cap{\mathcal C}^1([0,T], \mathop{\rm dom}B_1)\cap{\mathcal C}([0,T], \mathop{\rm dom}B_0) $,
удовлетворяющую уравнению \eqref{tsyp:2} и условиям \eqref{tsyp:1}.
\end{definition}

\smallskip

\begin{theorem}[4.1]Если пучок операторов $\vec{B}$ является $A$-диссипативным{\rm ,} то для любых $(v_{0},v_{1})\in \overline{\mathop{\rm im}(\mu L-M)^{-1}L}$ задача \eqref{tsyp:1}{\rm ,} \eqref{tsyp:2} имеет единственное решение  $v(t)${\rm .}
\end{theorem}

\smallskip

\begin{proof}
Сведём задачу \eqref{tsyp:1}, \eqref{tsyp:2} к задаче \eqref{tsyp:4}, \eqref{tsyp:5}, где
$u(t)=\begin{pmatrix} v(t)\\v'(t)\end{pmatrix}.$
Рассмотрим операторы $L=\begin{pmatrix}
I&0\\0&A\end{pmatrix}\in\mathcal L (\frak U;\frak F),$
$M=\begin{pmatrix} 0&I\\B_{0}&B_{1}\end{pmatrix}\in\mathcal Cl
(\frak U;\frak F),$ пространства $\frak U=\mathcal V\times
\mathcal V,$ $\frak F=\mathcal V\times \mathcal G.$

В силу леммы 2.3 пучок операторов $\vec{B}$ является $A$-диссипативным
точно тогда, когда оператор $M$ является $L$-диссипативным. Из
$L$-диссипативности оператора $M$ следует его
$(L, 0)$-радиальность.

По теореме 1.1  $\frak U^{1}= \overline{\mathop{\rm im}(\mu L-M)^{-1}L}$ является фазовым пространством уравнения \eqref{tsyp:4}. В силу определения 1.4 для
любых $u_{0}\in\frak U^{1}$ задача Коши для уравнения \eqref{tsyp:4} имеет
единственное решение $u(t)$. 
Значит задача \eqref{tsyp:1}, \eqref{tsyp:2} также однозначно разрешима при любых $\begin{pmatrix} v_{0}\\v_{1}\end{pmatrix}\in\frak U^{1}$, причём её решение $v(t)$ является первой компонентой вектора $u(t)$.\end{proof}

%% Благодарности, гранты и прочее

\begin{thebibliography}{9}

\RBibitem{zamtsyp:bib:Zamyshlyaeva} \by Замышляева~А.\,А. \paper
Фазовые пространства одного класса линейных уравнений соболевского
типа второго порядка \jour Вычисл. технол. \yr 2003
\vol 8
\issue 4
\pages 45--54
\misctransl
\by Zamyshlyaeva~A.\,A.
\paper Phase spaces of some class of linear second-order Sobolev-type equations
\jour Vychisl. Tekhnol. 
\vol 8
\issue 4
\pages 45--54 
\yr 2003

\Bibitem{zamtsyp:bib:Favini} \by Favini~A., Yagi~A. 
\book Degenerate differential equations is Banach spase 
\publaddr New York, NY 
\publ Marcel Dekker 
\yr 1999
\totalpages 312
\serial Pure and Applied Mathematics, Marcel Dekker
\vol 215


\Bibitem{zamtsyp:bib:Showalter} 
\by Showalter~R.\,E. 
\book  Monotone operators in Banach space and nonlinear partial differential equations
\serial Mathematical Surveys and Monographs
\vol 49
\publ American Mathematical Society
\publaddr Providence, RI
\yr 1997
\totalpages 278 

\Bibitem{zamtsyp:bib:Sviridyuk} 
\by  Sviridyuk~G.\,A., Fedorov~V.\,E. 
\book Linear Sobolev type equations and degenerate semigroups of operators
\serial Inverse and Ill-Posed Problems Series
\publaddr Utrecht
\publ VSP
\vol viii
\totalpages 216

\RBibitem{zamtsyp:bib:Fedorov} \by Федоров~В.\,Е. 
\paper Сжимающие полугруппы уравнений соболевского типа и относительно диссипативные операторы 
\jour Мат. замет. ЯГУ 
\yr 2001 
\issue 2
\vol 8 \pages 75--83
\misctransl
\by Fedorov~V.\,E.
\paper Contracting semigroups of Sobolev-type equations and relatively dissipative operators
\jour Mat. Zamet. YaGU 
\vol 8
\issue 2
\pages 75--83 
\yr 2001

\RBibitem{zamtsyp:bib:Fedorov1} \by Федоров~В.\,Е. 
\paper Обобщение теоремы Хилле--Иосиды на случай вырожденных полугрупп в~локально выпуклых пространствах
\jour Сиб. матем. журн.
\yr 2005
\vol 46
\issue 2
\pages 426--448
\mathnet{http://mi.mathnet.ru/smj977}
\mathscinet{http://www.ams.org/mathscinet-getitem?mr=2141208}
\zmath{http://www.zentralblatt-math.org/zmath/search/?an=Zbl 1106.47035}
\transl англ. пер.:~
\paper A generalization of the Hille--Yosida Theorem to the case of degenerate semigroups in locally convex spaces
\by Fedorov~V.\,E.
\jour Siberian Math. J.
\yr 2005
\vol 46
\issue 2
\pages 333--350
\crossref{http://dx.doi.org/10.1007/s11202-005-0035-9}

\end{thebibliography}



\noindent
\hskip6.5cm \thedate \par\noindent
\hskip6.5cm\therevdate




\makeengtitlem



 \newpage
%
% \tableofengcontents
%
%\end{document}
